\section{Aplicación web}
\label{sec:aplicacion-web}

\subsection{Arquitectura de la aplicación}

La aplicación web \texttt{apps/web} se construyó como una \emph{single
page application} (SPA) con React~18, TypeScript y Vite como
empaquetador. La organización del código sigue el patrón
\emph{feature-based}: cada dominio funcional (laboratorios, monitoreo,
reservas) vive en su propio directorio bajo \texttt{src/features/},
con sus componentes, llamadas a la API y pruebas colocalizadas, en
lugar de agruparse por tipo técnico. Los directorios
\texttt{src/router/}, \texttt{src/auth/}, \texttt{src/i18n/} y
\texttt{src/services/} concentran responsabilidades transversales
compartidas por todas las características.

\subsection{Enrutamiento y control de acceso}

El enrutamiento se implementó con \texttt{react-router-dom},
definiendo las cinco rutas principales exigidas por la guía, además de
las subrutas propias del módulo de reservas. La
Tabla~\ref{tab:rutas-web} resume las rutas y su protección.

\begin{table}[htbp]
\centering
\caption{Rutas de la aplicación web}
\label{tab:rutas-web}
\begin{tabular}{lll}
\toprule
Ruta & Descripción & Protegida \\
\midrule
/login & Inicio de sesión & No \\
/main & Panel principal & Sí \\
/settings & Configuración de usuario & Sí \\
/about & Información del sistema & No \\
/reservas & Listado de reservas & Sí \\
\bottomrule
\end{tabular}
\end{table}

Las rutas protegidas se envuelven mediante el componente
\texttt{ProtectedRoute.tsx}, que verifica la presencia de un token de
sesión válido antes de renderizar la ruta solicitada y redirige a
\texttt{/login} en caso contrario.

\subsection{Internacionalización y tema}

La aplicación soporta español e inglés mediante archivos de traducción
independientes bajo \texttt{src/i18n/locales/}. La preferencia de
idioma y el tema visual (claro u oscuro) se gestionan de forma
centralizada en \texttt{store/preferencesStore.ts} y persisten entre
sesiones. El comportamiento de ambos mecanismos está cubierto por
pruebas automatizadas (\texttt{settings-theme-i18n.spec.ts}).

\subsection{Autenticación y gestión de sesión}

El token JWT y su información asociada se almacenan en
\texttt{sessionStorage} en lugar de \texttt{localStorage}, decisión
deliberada para limitar la persistencia del token al ciclo de vida de
la pestaña del navegador y reducir la superficie de exposición ante
ataques de tipo \emph{cross-site scripting}. El contexto
\texttt{AuthContext.tsx} centraliza el inicio de sesión, la renovación
de token y el cierre de sesión, exponiéndolos al resto de la
aplicación mediante un \emph{hook} de React.

\subsection{Panel de monitoreo de laboratorios}

El panel de monitoreo, implementado en
\texttt{LaboratoriosPanel.tsx} y \texttt{MonitoreoPanel.tsx}, consume
el endpoint de detalle completo expuesto por el patrón Facade del
backend (\texttt{LaboratorioDetalleFacade}, Sección~\ref{sec:pruebas-cicd})
y la serie histórica de ocupación expuesta por
\texttt{PrometheusQueryClient}. El diseño visual se inspiró en la
identidad institucional de la UTEQ: encabezado en verde institucional,
tarjetas blancas con un borde izquierdo y un indicador luminoso
coloreado según el estado del laboratorio, una barra de capacidad
ocupada y tipografía monoespaciada (\texttt{JetBrains Mono}) para los
códigos de laboratorio. El cliente tipado \texttt{academicoApi.ts}
refleja exactamente los campos de \texttt{LaboratorioResponse.java}
del backend, evitando divergencias de contrato entre frontend y
backend. El componente cuenta con pruebas unitarias que cubren los
estados de carga, éxito, error y lista vacía.

\begin{figure}[htbp]
\centering
\includegraphics[width=\linewidth]{figuras/panel-monitoreo.png}
\caption{Panel de monitoreo de laboratorios}
\label{fig:panel-monitoreo}
\end{figure}

\subsection{Módulo de reservas}

El módulo de reservas implementa el flujo institucional desde el
\emph{gateway}, sin conexiones directas desde el navegador hacia los
microservicios internos. La autorización visual reutiliza los permisos
incluidos en la sesión: \texttt{SOLICITUD\_CREAR} habilita la nueva
solicitud, mientras \texttt{SOLICITUD\_LEER},
\texttt{RESERVA\_LEER} y \texttt{AGENDA\_GESTIONAR} controlan los
listados y el calendario. Estas comprobaciones mejoran la experiencia de
usuario, pero la decisión definitiva de acceso permanece en el backend.

\texttt{NuevaSolicitudPage} elimina la introducción manual de UUID. Para
un docente, la página resuelve primero su registro institucional a partir
del \texttt{perfilId} autenticado y carga mediante rutas canónicas las
materias, el período lectivo actual, los laboratorios y sus horarios. El
usuario selecciona opciones mediante nombres y códigos legibles, comprueba
la disponibilidad del laboratorio y envía la solicitud con una
\texttt{Idempotency-Key} estable durante los reintentos del mismo intento
lógico. La identidad del solicitante vuelve a ser validada y derivada por
el servicio de reservas.

La vista de reservas agrupa solicitudes y reservas respetando el ámbito
devuelto por el backend. Un docente observa sus propios registros; un
\texttt{ADMINISTRADOR\_PISO}, únicamente los laboratorios adscritos a su
piso; y \texttt{ADMINISTRADOR} conserva el alcance global. El detalle
muestra estado, fecha, horario, laboratorio, motivo, observación e
historial. Cuando existe una propuesta alternativa, el docente propietario
puede aceptarla o rechazarla, mientras el gestor autorizado puede poner la
solicitud en revisión, proponer otra franja, aprobarla o rechazarla.

La cancelación solo se presenta para estados admitidos por la máquina de
estados. En una solicitud aprobada, el backend cancela transaccionalmente
la solicitud y su reserva programada, evitando estados divergentes. El
calendario reutiliza el mismo alcance institucional y representa fecha,
hora, laboratorio y estado sin reconstruir reglas de seguridad en el
cliente.

Todas las llamadas utilizan un cliente HTTP centralizado con Bearer. Ante
un \texttt{401} se intenta una única renovación y se reintenta una vez la
petición original; un \texttt{403} no activa la renovación. Si la sesión
no puede renovarse, se limpian las credenciales de sesión y se redirige al
inicio de sesión. Los conflictos de estado o disponibilidad se presentan
como mensajes funcionales y no como excepciones internas.

\subsection{Empaquetado y despliegue}

La aplicación se empaqueta mediante un \texttt{Dockerfile} de
construcción en dos etapas: una primera etapa instala dependencias y
genera el artefacto de
